% An unstructured mesh has two kinds of element: the cells themselves and the
% faces between them. If each kind needs its own global reduction, then each
% kind wants its own communicator. (The original draws this with 3D cells; a
% 2D row carries the same point and stays legible at slide size.)
\documentclass[dvisvgm,border=3pt]{standalone}
\input{preamble.tex}
\begin{document}
\begin{tikzpicture}[x=1cm,y=1cm,
    cell/.style={draw=blu,line width=0.9pt,minimum size=9mm,
                 text=fg,font=\sffamily\small},
    face/.style={draw=uibkorange!70!black,line width=0.6pt,fill=uibkorange,
                 text=inkdark,font=\sffamily\scriptsize},
    note/.style={text=muted,font=\sffamily\scriptsize}]

  % --- the mesh: cells separated by faces ---------------------------------
  \foreach \l [count=\i from 0] in {A,B,C,D,E}{
    \node[cell] at ({\i*1.32},0) {\l};}
  \foreach \l [count=\i from 0] in {1,2,3,4}{
    \node[face,minimum width=3.2mm,minimum height=9mm] at ({\i*1.32+0.66},0) {\l};}

  % --- the two element types, each of which wants its own communicator ----
  \foreach \i in {0,...,3}{
    \node[face,minimum width=3.2mm,minimum height=9mm] at ({\i*0.42},-1.85) {};}
  \node[note,anchor=west] at (1.55,-1.85) {faces};

  \foreach \i in {0,...,4}{
    \node[cell,minimum size=7mm] at ({\i*0.86+3.60},-1.85) {};}
  \node[note,anchor=west] at (7.95,-1.85) {cells};

\end{tikzpicture}
\end{document}
